\subsection{Effective coherent local bases and higher-rank valuation examples}
\label{C:coherent-henselizations-subsection}

\begin{lemma}[Proved: finite henselian algebras without Noetherianity]
\label{C:nonnoetherian-henselian-lifting}
Let $(R,\mathfrak m)$ be any henselian local ring
with $2\in R^\times$.  Let $E$ be an associative
unital algebra finite as an $R$-module, with central
$R$-action, and put $J=\operatorname{rad}E$.
Then $E$ is semiperfect and idempotents lift through
every two-sided ideal contained in $J$.
For an $R$-linear involution $*$ and a self-adjoint
$a$ with $a-1\in J$, there exists a self-adjoint
$b\equiv1\pmod J$ with $b^2=a$.
Consequently nonsingular hermitian forms of either
sign on a fixed finite projective $E$-module are
classified up to isometry by reduction modulo $J$.
The same semisimple-block form classification
applies to any category of finite free rational
representations whose actual endomorphism rings
are finite over $R$.
\end{lemma}
\begin{proof}
We give the finite-algebra argument without a
Noetherian or completeness assumption.
For $h\in\mathfrak mE$ and $u\in E$, set
$c=1-uh$.  Left and right multiplication by
$c$ both reduce to the identity on the finite
$R$-module $E/\mathfrak mE$.
Nakayama applied to their cokernels makes
both maps surjective.  Thus $cv=1$ and
$wc=1$ for some $v,w\in E$.
The equality $w=w(cv)=(wc)v=v$ gives a
two-sided inverse.  The Jacobson criterion
proves $\mathfrak mE\subseteq J$.
Since $E/\mathfrak mE$ is finite dimensional
over the residue field, $E/J$ is semisimple
Artinian.  This use of two surjections
requires no assertion that arbitrary
surjective endomorphisms are injective.

If $A\subseteq E$ is a finite commutative
$R$-subalgebra, $E$ is a finite faithful
left $A$-module.  If $c\in A$ is invertible
in $E$ but belongs to a maximal ideal
$\mathfrak n$ of $A$, then
$E_{\mathfrak n}$ is a nonzero finite
$A_{\mathfrak n}$-module and
$cE_{\mathfrak n}=E_{\mathfrak n}$.
Nakayama is a contradiction.  Thus units
in $E$ belonging to $A$ are units in $A$,
and the Jacobson criterion gives
\[
                 A\cap J\subseteq\operatorname{rad}A.     \tag{C.3.NH1}
\]

Let $K\subseteq J$ and $v^2-v\in K$.
The determinant trick for multiplication
by $v$ on the finite $R$-module $E$
makes $A=R[v]$ finite over $R$.
A finite commutative algebra over a
henselian local ring is a finite product
of henselian local rings, without
Noetherianity
\cite[Lemma 10.153.4(1), Tag 04GH]
 {C-StacksCoherentHenselizations}.
By (C.3.NH1), $K\cap A$ lies in
the radical of each such local factor.
Every factor of $A/(K\cap A)$ is
local and has only the idempotents
$0$ and $1$.  Choose the same zeros
and ones in $A$.  The resulting
idempotent is congruent to $v$
modulo $K\cap A$, and hence modulo
$K$ in $E$.
Successive lifting in complementary
corners gives a complete set of
primitive orthogonal idempotents;
$E$ is semiperfect.

For the square root use $A=R[a]$.
The involution fixes this commutative
finite subalgebra pointwise, and
$a-1\in\operatorname{rad}A$.
In every henselian local factor
the polynomial $T^2-a$ has simple
residual root $1$.
The simple-root characterization
\cite[Lemma 10.153.3(2), Tag 04GG]
 {C-StacksCoherentHenselizations}
gives $b\in A$ with $b^2=a$
and $b\equiv1\pmod{\operatorname{rad}A}$.
Now $b+1$ is a unit of $A$, so
$b-1=(a-1)(b+1)^{-1}\in J$.
Also $b=b^*$ and $b$ is a unit.

All matrix and corner algebras
over $E$ are finite over $R$,
so the same argument applies
to endomorphisms of a fixed
finite projective $E$-module.
Residue form coefficients lift
and can be symmetrized using
$1/2$; invertibility modulo
the radical is invertibility.
For injectivity write two
forms as $q$ and $q a$.
Lift a residue isometry to
a unit $x$, so that
$d=x^*ax\equiv1\pmod J$.
The self-adjoint square root
$b^2=d$ gives
$(xb^{-1})^*a(xb^{-1})=1$.
Thus every residue isometry
lifts to an actual one.

In the representation category,
finite endomorphism rings are
semiperfect by the proved
algebra assertion.  Their
idempotents split as finite
projective, hence free,
$R$-summands.  The finite
Krull--Schmidt exchange
argument gives unique
indecomposable decompositions.
The radical quotient of
the endomorphism category
has the division-algebra
blocks of those summands.
Duality exchanges paired
blocks and acts by adjoints
on self-dual blocks.
The preceding exact form
lifting identifies its
multiplicity-form classification
with the actual one, exactly
as in the transfer argument
of Lemma~\ref{C:henselian-finite-algebra-lifting}.
All rings to which that
argument is applied here
are finite endomorphism
matrix or corner algebras;
the lifting proof above
supplies its hypotheses.
\end{proof}

\begin{theorem}[Proved: henselization algorithms from effective coherent local data]
\label{C:effective-coherent-henselization-WG}
Let $B$ be a countable commutative local ring
with exact finite names, decidable name validity
and equality, computable arithmetic, and a
computable residue map onto a field $\kappa$
of characteristic $p>2$.
Assume an ordinary terminating algorithm
computes a finite generating set for the
kernel of every finite matrix over $B$.
For $\kappa$ assume either a supplied finite
field table, or the effective infinite
$C_1$ field data preceding
Lemma~\ref{C:C1-henselization-access}.
Put $R=B^h$.

Then $R$ has exact finite names, computable
arithmetic, equality, residues, unit inverses
and residue lifts, and full finite matrix-kernel
algorithms.  All finite projective rational
$O_{m,R}$-indecomposables have a decidable
literal canonical index set.
Module isomorphism, indecomposability, duality,
form sign, actual isometry and split-group
equality are decidable by ordinary terminating
algorithms.
The canonical additive description is
(C.3.HD2) for finite residue and (C.3.CH3)
for infinite $C_1$ residue, with all their
actual generators, multiplication and every
specified exterior-power coefficient computed
in finite time.
No Noetherianity, separated adic topology,
or presentation by flat Noetherian stages
is assumed on $B$ or $R$.
\end{theorem}
\begin{proof}
The kernel hypothesis implies in particular
that $B$ is coherent, but we use the actual
algorithm, not just this existence property.
Since $B$ is local, an element is a unit
exactly when its residue is nonzero.
If needed its inverse is found by enumerating
$B$ and testing the equation $uv=1$;
existence proves termination.
Likewise every residue element has a lift
found by enumeration and its decidable
residue equality test.

Raw names in $R$ use a monic $f\in B[U]$,
a simple residue root $a\in\kappa$,
a polynomial $h$ with
$\overline{f'}(a)\overline h(a)\ne0$,
and an element in
\[
 C_g,\qquad C=B[U]/(f),\qquad g=f'h.                     \tag{C.3.NH2}
\]
The algebra $C$ is finite free over $B$
in its power basis.  Its multiplication
table is computed by monic division.
The localization is a pointed standard
\'etale neighborhood and maps to $R$
by the chosen root.
All data have decidable validity.
All elements of $R$ occur in these names
by the pointed \'etale construction and
standard local structure of \'etale maps
\cite[Lemma 10.155.1 and Proposition 10.144.4,
Tags 0BSK and 00UE]{C-StacksCoherentHenselizations}.
No approximation by powers of the maximal
ideal is being used.

Finitely many such neighborhoods may be
combined by tensoring their finite free
algebras over $B$ and inverting the product
of the corresponding $g$'s.
The tensor algebra is again explicitly
finite free over $B$.  Denote it also by
$C$, the inverted element by $g$, and
the resulting residue map by
$\rho:C_g\to\kappa$.
It sends each root variable to its
specified residue root and is surjective.
If $\mathfrak q=\ker\rho$, then
$(C_g)_{\mathfrak q}^h=R$.
Hence $C\to R$ is flat and
$(C_g)_{\mathfrak q}\to R$ is
faithfully flat
\cite[Section 10.155 and Lemma 15.46.1,
Tags 0BSK and 07QL]{C-StacksCoherentHenselizations}.

For $z=c/g^n$, multiplication by $c$
is a finite $B$-matrix on the free
$B$-module $C$.
The supplied algorithm computes
$B$-generators of its kernel
$\operatorname{Ann}_C(c)$; write them as
$u_1,\ldots,u_s$.
They also generate that ideal as
a $C$-module: every element already
is a $B$-linear combination of them.
Exactness of localization gives
$\operatorname{Ann}_{C_g}(z)
 =\operatorname{Ann}_C(c)_g$.
Faithful flatness of the pointed
local henselization now gives
\[
 z=0\text{ in }R
 \quad\Longleftrightarrow\quad
 \rho(u_j)\ne0\text{ for some }1\leq j\leq s.             \tag{C.3.NH3}
\]
Thus equality is decided by a finite
test on the computed generators.
This argument also applies to a
difference of names in their tensor
neighborhood.

Arithmetic takes place in that
neighborhood.  A unit numerator
has nonzero residue, so adjoining
its inverse is another pointed
\'etale localization and gives
an exact inverse name.
The given base residue lifts
serve also as lifts in $R$.
Internal names can be returned
to the standard list by an
equality search, whose termination
follows from coverage.
Retaining each valid natural-number
code only when no smaller valid
code is equal to it gives canonical
finite names by a finite test.

For a finite matrix with entries
in $C_g$, multiply it by a common
power of the unit $g$ to obtain
a matrix $A_C:C^u\to C^v$.
Regard this as a finite $B$-matrix
between free $B$-modules.
Its computed $B$-kernel generators
also generate the full $C$-kernel.
Flatness of $C\to R$ identifies
their base change with the entire
actual kernel over $R$.
Multiplication by the cleared
unit denominator does not change
that kernel.  This proves all
claimed exact access to $R$.

In particular every actual
representation endomorphism ring
is finite as an $R$-module:
it is the kernel of its finite
coefficient-commutator system.
Apply Lemma~\ref{C:nonnoetherian-henselian-lifting}
to those endomorphism rings
and their finite matrix and
corner algebras.
It supplies semiperfectness,
Krull--Schmidt decomposition
and exact radical isometry
lifting without Noetherianity.

For finite $\kappa$, every
step of the full proof of
Theorem~\ref{C:effective-henselian-DVR-WG}
now applies with the computed
kernel generators in place
of its Smith bases.
Specifically the actual Hom
residue images are finite
vector spaces, so invertible
intertwiners and form signs
are decided by enumeration.
Their finite endomorphism
image algebras have computable
radicals and semisimple
quotients.  A positive
nontrivial-idempotent decision
starts an enumeration of
actual lifts, which terminates
by the preceding lemma.
Independent residue columns
split the free summands;
rank decreases at each step.
The same least-code index
tests and finite residue-form
classification give exactly
(C.3.HD2), including its
ordinary isometry decisions
and actual generators.

For infinite effective $C_1$
residue use the proof of
Theorem~\ref{C:C1-henselization-WG}.
Its finite determinant grids,
positive finite-tower certificates
for the actual residue
endomorphism algebras, least
finite module codes, square-class
factorization tests and
residual Gram extraction
use precisely the finite
kernel access and exact
henselian lifting just proved.
They do not require any
additional Noetherian step.
They give all of (C.3.CH3)
and its exact actual isometry
and equality decisions.

In both cases tensor products
and exterior minors are finite
actual matrix constructions
followed by the proved
classification procedure.
Finite virtual coefficients
use bilinearity and inversion
of $\lambda_t$ to the requested
order.  All searches invoked
in these algorithms are the
positive searches with proved
existence just specified.
This proves every claimed
ordinary terminating output.
\end{proof}

\begin{corollary}[Proved: effective valuation domains and explicit nonseparated families]
\label{C:lexicographic-valuation-henselization-WG}
Let $V$ be a countable valuation domain with
exact finite names and decidable arithmetic
and equality, decidable divisibility and
computable division whenever one element
divides another.  Suppose its residue
map and field have the effective data
required in
Theorem~\ref{C:effective-coherent-henselization-WG}.
Then that theorem gives complete ordinary
canonical classification, equality and
all split $WG$ operations over $V^h$.

Here are actual uniform families supplying
all this access.  For $r\geq2$, take either
\[
 \begin{array}{ll}
 K=\mathbf F_q(x_1,\ldots,x_r),
     &v(K^\times)=\mathbf Z^r_{\rm lex},\quad
       v(x_i)=e_i,\\
 K=\mathbf Q(x_1,\ldots,x_{r-1}),
     &v(K^\times)=\mathbf Z^r_{\rm lex},\quad
       v(x_i)=e_i,\quad v(p)=e_r,
 \end{array}                                              \tag{C.3.NH4}
\]
where $q$ is an odd prime power and $p$ an
odd prime.  In the first line the valuation
of a polynomial is its least exponent
vector in lexicographic order.
In the second line it is the least
exponent vector in the variables,
followed by the $p$-adic valuation
of that term's nonzero rational
coefficient.  Extend to rational
functions by subtraction, and put
$V=\{0\}\cup\{f\in K^\times:v(f)\geq0\}$.

These $V^h$ are countably infinite
non-Noetherian henselian local domains,
with residue field $\mathbf F_q$
or $\mathbf F_p$, respectively.
Their maximal ideals are generated
by $y=x_r$ in the first family
and $y=p$ in the second, yet
\[
             0\ne x_1\in\bigcap_{n\geq0} y^nV^h.       \tag{C.3.NH5}
\]
They are therefore not separated for
their maximal-ideal-adic topology,
and cannot be filtered colimits
of Noetherian local rings along
flat local transition maps.
In particular these complete ordinary
answers are outside the earlier
flat-Noetherian-stage families.
\end{corollary}
\begin{proof}
Over the supplied valuation domain,
finite matrix kernels are computable.
Choose a nonzero entry dividing all
other entries, found by the decidable
total divisibility relation on a
valuation domain.  Permute it into
the first position and clear its
column and row using the allowed
exact divisions and elementary
invertible operations.
Repeat on the remaining smaller
matrix.  This diagonalization
terminates after finitely many
pivots.  Since $V$ is a domain,
the nonzero diagonal entries
have zero kernels; the remaining
coordinates give a free kernel
basis, transported back by the
computed invertible column
operations.  Thus every hypothesis
on $B=V$ in the theorem is supplied.

We verify the explicit families.
In the first line of (C.3.NH4),
the least term of a product is
the product of the two least
terms.  In the second line the
same is true of the leading
variable exponent, and its
rational coefficient has the
additive $p$-adic valuation.
Thus in both lines
$v(fg)=v(f)+v(g)$ and
$v(f+g)\geq\min\{v(f),v(g)\}$.
The quotient definition is
independent of the chosen
rational-function presentation.
Every value in $\mathbf Z^r$
is attained by a Laurent
monomial, using $p$ as the
last variable in the second
line.  The set $V$ is a
valuation domain because,
for any nonzero rational
function, either it or its
inverse has nonnegative value.

All arithmetic, equality
and valuations are computed
from finite rational-function
data.  Divisibility in $V$
is the test $v(a)\leq v(b)$
for $a\mid b$, and its
quotient is the rational
function $b/a$.
The residue of a positive
value is zero.
For value zero, in the
first family it is the
ratio of the leading
coefficients in $\mathbf F_q$;
in the second it is that
rational ratio, a $p$-adic
unit, reduced modulo $p$.
These rules are independent
of presentation and respect
addition and multiplication
by the valuation inequalities.
Their kernels are exactly
the positive-value elements,
and constants supply all
residues.  Thus the residue
fields and all required
maps and lifts are explicit.

The least positive value
in $\mathbf Z^r_{\rm lex}$
is $e_r$.  Every element
of the maximal ideal is
therefore divisible by
the indicated $y$, so
that ideal is $yV$.
The element $x_1/y^n$
has value $e_1-ne_r>0$
for every $n$; hence it
belongs to $V$.
Henselization is faithfully
flat with the maximal ideal
extended from the source
\cite[Lemma 15.46.1, Tag 07QL]
 {C-StacksCoherentHenselizations}.
Thus $yV^h$ is the maximal
ideal and (C.3.NH5) holds
with $x_1\ne0$.

For completeness $V$ is
normal.  If $z\in K$ had
negative value and satisfied
a monic equation over $V$,
the leading power of $z$
would have strictly smaller
value than every other
term, an impossibility
for a zero sum.
Henselization of a normal
local domain is a normal
domain even without
Noetherianity
\cite[Lemma 15.46.6, Tag 07QL]
 {C-StacksCoherentHenselizations}.
Thus $V^h$ is a domain.
It is countable by its
finite names and infinite
because it contains $V$.
Equation (C.3.NH5) and
Krull's intersection theorem
\cite[Lemma 10.51.4, Tag 00IP]
 {C-StacksCoherentHenselizations}
show that it is not
Noetherian.

Suppose it were a filtered
colimit of Noetherian local
rings with flat local
transition maps.
The maps from every stage
to the local colimit are
flat and local, hence
faithfully flat and
injective.  Choose a
stage $A$ containing
the two elements $x_1,y$.
The element $y$ lies
in its maximal ideal,
since the stage map
is local.
For every $n$, faithful
flatness reflects ideal
membership, so
$x_1\in y^nV^h$ implies
$x_1\in y^nA$.
Krull's intersection
theorem in $A$ gives
$x_1=0$, contradicting
injectivity and (C.3.NH5).
This proves the strict
scope separation and all
claims in the corollary.
\end{proof}
